\subsection{联立求解：同时拟合两条记忆}
直接累加外积时，每条新记忆可能改变旧 Key 的输出。
可以让矩阵同时拟合全部输入与目标：
\[
\boxed{\min_{\mathcal{M}}\sum_{i=1}^{T_k}\|k_i\mathcal{M}-v_i\|^2}.
\]
这叫最小二乘：让同一个状态同时满足各条关联。

刚才的外积相加得到
\[
\mathcal{M}_{\rm simple}=\begin{bmatrix}1&0.6\\0&0.8\end{bmatrix}.
\]
要同时读对 A、B，可以直接联立求解。设矩阵的两行为 $m_1,m_2$。
A 要求 $m_1=(1,0)$；B 要求 $0.6m_1+0.8m_2=(0,1)$，所以
\[
m_2=\frac{(0,1)-0.6(1,0)}{0.8}=(-0.75,1.25).
\]
于是
\[
\boxed{\mathcal{M}_{\rm better}=\begin{bmatrix}1&0\\-0.75&1.25\end{bmatrix}},
\]
就能逐项验算：
\[
(1,0)\mathcal{M}_{\rm better}=(1,0),\qquad
(0.6,0.8)\mathcal{M}_{\rm better}=(0,1).
\]
负权重在这里用于\textbf{抵消交叉干扰}。
但如果约束不相容，线性层并不能任意精确存取无限多条关联。

\subsection{从误差修正到 Delta 写入}
\phantomsection\label{concept:delta-rule}
比无条件相加外积更进一步：写入新配对之前，先看旧状态
$\mathcal{M}_{t-1}$ 已经读出什么，只写入尚未记住的残差：
\[
r_t=v_t-k_t\mathcal{M}_{t-1},\qquad
\boxed{\mathcal{M}_t=\mathcal{M}_{t-1}+\eta k_t^\top r_t}.
\]
这里 $k_t:1\times d_k$、$r_t:1\times d_v$；
若 $k_t$ 归一化且 $\eta=1$，可以在该 Key 方向消除\textbf{当次}读出误差。
但这未必保护此前已学会的 Key。

\subsection{［复算］写好 B 时，A 会不会再次出错？}

仍用两条单位 Key 和 Value：
\[
k_A=(1,0),\ v_A=(1,0),\qquad
k_B=(0.6,0.8),\ v_B=(0,1).
\]
从 $\mathcal{M}_0=0$、$\eta=1$ 开始，先写 A：
\[
\mathcal{M}_1=k_A^\top v_A=\begin{bmatrix}1&0\\0&0\end{bmatrix}.
\]
写 B 前，旧状态的读取为 $k_B\mathcal{M}_1=(0.6,0)$，
所以不是直接追加 $k_B^\top v_B$，而是计算残差：
\[
r_B=v_B-k_B\mathcal{M}_1=(-0.6,1).
\]
下一步得到
\begin{equation}
\begin{aligned}
\mathcal{M}_2&=\mathcal{M}_1+k_B^\top r_B\\
&=\begin{bmatrix}1&0\\0&0\end{bmatrix}
+\begin{bmatrix}0.6\\0.8\end{bmatrix}(-0.6,1)
=\boxed{\begin{bmatrix}0.64&0.6\\-0.48&0.8\end{bmatrix}}.
\end{aligned}
\label{eq:delta-overwrites-key-a}
\end{equation}
立即查询 B 得到\textbf{正确的} $k_B\mathcal{M}_2=(0,1)$；
反过来，旧的 A 查询却变成
$k_A\mathcal{M}_2=(0.64,0.6)$，不再等于 $(1,0)$。
A 的误差范数为
\[
\|k_A\mathcal{M}_2-v_A\|_2=\sqrt{(-0.36)^2+0.6^2}
\approx\boxed{0.700}.
\]
若直接相加外积，A 的误差为 $0.600$。
\textbf{在这个明确例子里，单步 Delta 对旧记忆的扰动反而更大。}

\begin{boundarybox}
Delta 写入的价值是\textbf{按残差修正当前读取}，不意味着
每写一条后所有旧记忆都更准确。
交替重访 A、B，在这两条线性无关的单位 Key 与 $\eta=1$
的条件下可收敛到本章的精确解 $\mathcal{M}_{\rm better}$。
它是一次次更新得到的结果，而不是两次写入后的立即结果。
\end{boundarybox}


\begin{idea}
最小二乘是一个全局优化目标；Delta 写入是一种依赖当前状态的局部更新。
两者都在减少干扰：前者一次求联立解，后者逐条修正当前读出。
\hyperref[sec:modern-delta-rule]{现代架构章}用同一对 Key 看到 B 被精确修正；
这里补上另一半：旧 Key A 随之被扰动；Gated DeltaNet 的遗忘门控就在调节旧记忆保留多少。
\end{idea}

\noindent\emph{对应迁移练习：}\hyperref[exercise:delta-reverse]{附录 C 的 C-107}
交换 A、B 的写入次序，检查被扰动的是否始终是旧 Key。

\chapterchecktasks{04-key-value-outer-product}{理解、手算与判断}{
\label{check:06-better-linear-memory}
\item\textbf{形状：}行向量 Key $k:1\times2$、Value $v:1\times3$。外积写入 $\mathcal{M}=k^\top v$ 的形状是什么？查询 $q:1\times2$ 后 $q\mathcal{M}$ 又是什么形状？\chapterinlinehint{先写转置 Key 的形状，再核对内侧维度。}
\item\textbf{预测：}同一条记忆若把 Key 写入时改为 $k'=(0.5,0)$，再以 $q=k'$ 读取，恢复的 Value 是多少倍？\chapterinlinehint{把写入和读取两次出现的 Key 缩放都算上。}
\item\textbf{手算：}$k_A=(1,0),k_B=(0.6,0.8)$；$v_A=(1,0),v_B=(0,1)$。把两条外积相加，用 $k_A$ 查询得到什么？\chapterinlinehint{分别求两条外积，再相加并用查询读取。}
\item\textbf{预测：}若把 $k_B$ 改为 $(0,1)$，保持其他项不变，$k_A$ 查询时串扰如何变化？\chapterinlinehint{先比较修改前后的交叉点积。}
\item\textbf{定位：}Delta 写入为什么先用当前状态读取一次？“期望 Value 减去已有读出”代表什么量？\chapterinlinehint{先写目标读出与当前读出的差。}
\item\textbf{手算：}写入 B 后 $\mathcal{M}_2=\begin{bmatrix}0.64&0.6\\-0.48&0.8\end{bmatrix}$；用 $k_A=(1,0)$ 读出什么？与目标 $v_A=(1,0)$ 相差多少？\chapterinlinehint{先用查询选出对应行，再减去目标向量。}
\item\textbf{找错：}同学说“B 的残差为零后，原先的 A 也必定被完全修复”。指出更新为何可能影响其他查询。\chapterinlinehint{考察新增外积与另一条 Key 的点积。}
}

\par\smallskip
减少串扰的线性办法都在\textbf{写入端}：联立求解或逐条按残差修正。
下一章转到\textbf{读取端}，让匹配分数先经过非线性筛选再读出 Value。
